\chapter*{Experiment Mode}
\section*{Conditions}
This file was generated to see the results of three different algorithms solving multiple 100 random generated 0/1 Knapsack Problem and 
store the diverse results everytime the algorithm solves each problem. The experiment done has specific conditions everytime the \verb|knapsack| 
program is run.

\begin{itemize}
    \item The experiment tests 100 different 0/1 Knapsack Problems.
    \item When running the experiment mode with the \verb|knapsack| program, it needs to receive a parameter of how many iterations to do. This
    is the parameter \verb|-E=n|, where \verb|n| indicates the amount of iterations to do the experiment.
    \item Each problem will be solved by three algorithms and data will be collected and stored in tables with the same format 
    as shown in Table \ref{tab:table_format}.
    \item Each cell represents a specific problem. The rows indicate the size of the knapsack for the the problem and the columns
    indicate how many items are in the problem.
    \item Each problem also has specific conditions, because each iteration of the table generates a random problem:
        \begin{itemize}
            \item Value of items are in the range of 0 and 100.
            \item Cost of items are in the range of 0 and 40\% of the knapsack capacity.
        \end{itemize}
    \item When running the program with the amount of iterations, it will generate 100 random problems with the conditions explained before,
    it will solve each other, store data and it will do other 100 random problems until the amount of iterations are satisfy.
    \item The data collected for the tables will be an overall average or percentage of all iterations done.
    \item At the end the experiment returns 5 tables.
    \begin{enumerate}
        \item Table of average times needed to solve each problem with the Dynamic Programming Algorithm.
        \item Table of average times needed to solve each problem with the Basic Greedy Algorithm.
        \item Table of average times needed to solve each problem with the Proportional Greedy Algorithm.
        \item Table of percentages of how often the Basic Greedy Algorithm finds the optimal solution of a problem.
        \item Table of percentages of how often the Proportional Greedy Algorithm finds the optimal solution of a problem.
    \end{enumerate}
\end{itemize}

\begin{table}[h]
\centering
\begin{tabular}{|c|c|c|c|c|c|c|c|c|c|c|}
\hline
& \multicolumn{10}{c|}{Amount of items} \\
\hline
Knapsack & 10 & 20 & 30 & 40 & 50 & 60 & 70 & 80 & 90 & 100 \\
\hline
100 & & & & & & & & & & \\
\hline
200 & & & & & & & & & & \\
\hline
300 & & & & & & & & & & \\
\hline
400 & & & & & & & & & & \\
\hline
500 & & & & & & & & & & \\
\hline
600 & & & & & & & & & & \\
\hline
700 & & & & & & & & & & \\
\hline
800 & & & & & & & & & & \\
\hline
900 & & & & & & & & & & \\
\hline
1000 & & & & & & & & & & \\
\hline
\end{tabular}
\caption{Format of the tables}
\label{tab:table_format}
\end{table}

\newpage

\chapter*{Algorithms}

\section*{Dynamic Programming Algorithm}

This algorithm uses a table to get to the optimal solution for the generated 0/1 Knapsack Problem. The strategy here is to solve
subproblems, write down the optimal solution for that subproblem on a table and by using those solutions build the solution 
for the original problem. It always find the optimal solution for a 0/1 Knapsack Problem, as explained in course of Operations Research.

\section*{Basic Greedy Algorithm}

This algorithm focuses on trying to put inside the knapsack the most valuable item, if it does not fit in, tries with the next
most valueable item until the are no more items available.

\section*{Proportional Greedy Algorithm}

It has the same strategy as the Basic Greedy, but it focuses on trying to put inside items based by their proportion between its
value and cost. The item with highest proportion goes first, if it does not fit, goes to the next item with the highest proportion.

